\section{Τεχνικό Υπόβαθρο}

Στο κεφάλαιο αυτό, θα χτιστεί το τεχνικό υπόβαθρο που χρειάζεται για την κατανόηση της εργασίας. Συμπεριλαμβάνονται λεπτομέρειες τόσο για έννοιες μηχανικής μάθησης όσο και για βασικές ιδέες της γλωσσολογίας.

\subsection{Φωνήματα και Visemes}
\label{Phonemes}
Τα φωνήματα αποτελούν το βασικό αντικείμενο μελέτης της συγκεκριμένης εργασίας. Ένα φώνημα ορίζεται ως ένα σύνολο παρόμοιων ήχων ομιλίας το οποίο γίνεται αντιληπτό από ομιλητές μίας γλώσσας ως ένας βασικός ήχος, πρόκειται δηλαδή για τη μικρότερη πιθανή μονάδα φωνής. Συμβολίζονται συνήθως ως ένας χαρακτήρας εντός δύο forward slash. Για παράδειγμα, το φώνημα /b/ εκπροσωπεί τον ήχο στην αρχή της αγγλικής λέξης bat. Κάθε ομιλούμενη γλώσσα αξιοποιεί τα φωνήματα σε διαφορετικό βαθμό, ενώ ίδια γράμματα ή φθόγγοι μπορεί σε διαφορετικές γλώσσες να προφέρονται τελείως διαφορετικά. Τα προβλήματα αυτά προσπαθεί να λύσει το International Phonetic Alphabet (IPA), το οποίο είναι ένα ``αλφάβητο'' από φωνήματα βασισμένο κυρίως στους λατινικούς χαρακτήρες και ομοιόμορφο για κάθε γλώσσα. Να σημειωθεί ότι το πλήρες IPA είναι υπερβολικά περίπλοκο για την ανάλυση που ακολουθεί, καθώς αρκετά φωνήματα συναντώνται πολύ σπάνια (ή καθόλου) στην αγγλική γλώσσα που μελετάται καθ' όλη τη διάρκεια της εργασίας.

`Οντας η ελάχιστη μονάδα ομιλίας, τα φωνήματα έχουν αποτελέσει το αντικείμενο ποικίλων γλωσσολογικών μελετών, πολλές από τις οποίες αφορούν το ``διάβασμα χειλιών'' (lip reading) ή, αλλιώς, την κατανόηση ομιλίας χρησιμοποιώντας αποκλειστικά οπτικά (και όχι ακουστικά) μέσα. Από τις μελέτες αυτές έχει δημιουργηθεί το πρόβλημα των ``visemes'', τα οποία αποτελούν το σύνολο ήχων ομιλίας με πανομοιότυπη όψη στα χείλια (ή γενικότερα στο πρόσωπο). Να σημειωθεί ότι τα σύνολα αυτά δεν περιορίζονται σε μοναδικά φωνήματα: για παράδειγμα, οι φράσεις ``elephant juice'' και ``i love you'' αποτελούν ένα χιουμοριστικό αλλά έγκυρο παράδειγμα viseme. Στα πλαίσια της συγκεκριμένης μελέτης, στο εξής ως visemes θα αναφέρονται τα φωνήματα κατά την έκφανση των οποίων η όψη του ομιλητή είναι όμοια. Για παράδειγμα, τα φωνήματα /m/, /b/ και /p/ αποτελούν ένα παράδειγμα viseme. Τα visemes προσφέρουν πληροφορία για την τοποθεσία του στόματος στην οποία παράγεται ο ήχος, αλλά όχι για τον τρόπο με τον οποίο παράγεται ο ήχος, το οποίο απαιτεί ακουστικά στοιχεία.

Πέραν της ομαδοποίησης τους μέσω των visemes, τα φωνήματα μπορούν να ομαδοποιηθούν και με βάση 3 ακόμα χαρακτηριστικά: τον τόπο, τον τρόπο εκφοράς τους και το αν είναι έμφωνα ή όχι. Για να παρουσιαστούν με μεγαλύτερη ακρίβεια τα παραπάνω, συμπεριλαμβάνεται το επίσημο chart φωνημάτων του IPA. Αγνοώντας τις υπερβολικές λεπτομέρειες στο Σχήμα \ref{IPAChart} και εστιάζοντας την προσοχή στον πρώτο πίνακα με τίτλο CONSONANTS (PULMONIC) και στον αμέσως κάτω δεξιά χάρτη με τίτλο VOWELS, μπορούν να εξαχθούν τα παρακάτω συμπεράσματα.

Παρατηρούμε ότι τα σύμφωνα χωρίζονται σε 11 κατηγορίες όσον αφορά τον τόπο εκφοράς τους: τα δειχιλικά (bilabial, παράγονται χρησιμοποιώντας τα 2 χείλη), τα χειλοδοντικά (labiodental, παράγονται με τη συμμετοχή χειλιών και δοντιών), τα οδοντικά (dental, με τα δόντια), τα φατνιακά (alveolar, παράγονται με τη γλώσσα να ακουμπά ή να πλησιάζει τη φατνιακή ακρολοφία), τα μεταφατνιακά (postalveolar, παράγονται με τη γλώσσα να τοποθετείται πίσω από τη φατνιακή ακρολοφία), τα ανακεκαμμένα (retroflex, παράγονται από την άκρη της γλώσσας όταν κάμπτεται προς τα πίσω προς τον ουρανίσκο), τα ουρανικά (palatals, παράγονται υψώνοντας το κυρίως τμήμα της γλώσσας στον σκληρό ουρανίσκο), τα υπερωικά (velar, όταν το πίσω μέρος της γλώσσας αγγίζει ή πλησιάζει τον μαλακό ουρανίσκο), τα σταφυλικά (uvular, με το πίσω μέρος της γλώσσας να αγγίζει ή να πλησιάζει τη σταφυλή), τα φαρυγγικά (pharyngeal, παράγονται περιορίζοντας τον φάρυγγα) και τέλος τα γλωττιδικά (glottal, χρησιμοποιώντας τη γλωττίδα). Το χαρακτηριστικό αυτό είναι που επηρεάζει κατά κύριο λόγο τα visemes, καθώς ο τρόπος με τον οποίον παράγεται ένα φώνημα και το αν είναι έμφωνο ή όχι πολλές φορές γίνεται κρυφό λόγω της μορφολογίας των χειλιών, που επηρεάζεται προφανώς από τον τόπο εκφοράς του φωνήματος.

Πέραν του τόπου εκφοράς, τα σύμφωνα διαχωρίζονται στον πίνακα με βάση τον τρόπο εκφοράς τους ως εξής: τα έκκροτα/κλειστά σύμφωνα (plosive, παράγονται όταν διακόπτεται τελείως η ροή του αέρα και απελευθερώνεται ξαφνικά), τα ρινικά (nasal, επιτρέπουν στον αέρα να διαφύγει μέσω της ρινικής κοιλότητας και όχι του στόματος), τα μέσω τρίλιας (trill, παράγονται μέσω γρήγορης δόνησης λόγω της ροής αέρα), τα παλμού/κτύπου (tap/flap, όταν η γλώσσα ακουμπάει στιγμιαία ένα σημείο του στόματος), τα τριβόμενα (fricative, όταν ο αέρας περνάει από στενή δίοδο δημιουργώντας τριβή), τα πλευρικά τριβόμενα (lateral fricative, όπου ο αέρας παράγει τριβή καθώς περνάει από τα πλάγια της γλώσσας), τα προσεγγιστικά (approximant, όπου δύο αρθρωτές πλησιάζουν χωρίς να παράγουν τυρβώδη ροή αέρα ή ολικό κλείσιμο), και τα πλευρικά προσεγγιστικά (lateral approximant, μπλοκάροντας το κέντρο της στοματικής κοιλότητας στην περιοχή του αρθρωτή και επιτρέποντας τη ροή του αέρα από τα πλάγια). Πειραματιζόμενοι με τα παραπάνω μπορούμε να συνειδητοποιήσουμε πώς παράγονται τα φωνήματα που χρησιμοποιούμε, όπως και να αρθρώσουμε φωνήματα που ίσως δεν έχουμε συνηθίσει (π.χ. στα ελληνικά δεν υπάρχει πλευρικό τριβόμενο φώνημα, αλλά σε γλώσσες όπως τα ουαλικά υπάρχει).

Τέλος, ο χαρακτηρισμός έμφωνο ή άφωνο φαίνεται από τη θέση των συμφώνων στο εκάστοτε κελί. Αυτά που είναι τοποθετημένα προς τα δεξιά είναι έμφωνα, ενώ αυτά προς τα αριστερά είναι άφωνα. Τα φωνήεντα, από την άλλη είναι όλα έμφωνα και διαχωρίζονται με βάση το ύψος της γλώσσας (στον κατακόρυφο άξονα), την προώθησή της (στον οριζόντιο άξονα) και το αν τα χείλη είναι στρογγυλεμένα (rounded) ή όχι. Όσον αφορά το ύψος, οι κατηγορίες είναι close, close-mid, open-mid και open, ενώ στον οριζόντιο άξονα οι πιθανές επιλογές είναι front, central και back. Όπου στον χάρτη εμφανίζονται δύο φωνήματα, το δεξιότερο παράγεται με στρογγυλεμένα χείλη.

Με γνώση των παραπάνων και αντίστοιχη προσαρμογή στις ανάγκες της εκάστοτε γλώσσας ή, στην περίπτωση περιορισμένου περιβάλλοντος, στις ιδιαιτερότητες του λεξιλογίου, τα φωνήματα μπορούν να ομαδοποιηθούν με τους παραπάνω κανόνες. Έτσι, σε πλαίσια lip reading, είναι δυνατόν να χτιστούν μέθοδοι εφαρμογής διαφορετικής ποινής με βάση την ομοιότητα των φωνημάτων (είτε σε επίπεδο visemes είτε σε επίπεδο χαρακτηριστικών, όπως παρουσιάστηκε) ή και μέθοδοι πρόβλεψης λέξεων που να συνυπολογίζουν την ομοιότητα φωνημάτων κατά την αντικατάστασή τους (για περισσότερα βλ. παράγραφο \ref{PhonemeExtraction}).

\begin{figure}[h]
	\centering
	\includegraphics[width=\textwidth]{IPAChart.png}
	\caption{Επίσημο chart φωνημάτων του IPA, έκδοση 2020}
	\label{fig:IPAChart}
\end{figure}

